\documentclass[10pt, twocolumn]{article}
\usepackage[margin=0.75in, top=0.6in, bottom=0.6in]{geometry}
\usepackage{newtxtext}
\usepackage{titlesec}
\usepackage{enumitem}
\usepackage{xurl}
\usepackage{hyperref}
\usepackage{xcolor}

% Spacing
\setlist{nosep, leftmargin=*}
\linespread{1.05}
\setlength{\parskip}{3pt}
\hyphenpenalty=10000
\exhyphenpenalty=10000
\renewcommand{\footnotesize}{\scriptsize}

% Section formatting with visual distinction
\titleformat{\section}{\normalsize\bfseries\scshape}{\thesection}{0.5em}{}[\vspace{1pt}\titlerule]
\titlespacing*{\section}{0pt}{6pt}{3pt}
\titleformat{\subsection}{\small\bfseries\itshape}{\thesubsection}{0.5em}{}
\titlespacing*{\subsection}{0pt}{4pt}{2pt}

\hypersetup{
    colorlinks=true,
    linkcolor=black,
    urlcolor=blue,
    citecolor=black,
    breaklinks=true,
}

\pagestyle{plain}

% Compact Bibliography
\makeatletter
\renewenvironment{thebibliography}[1]
     {\section*{References}%
      \scriptsize
      \setlength{\topsep}{0pt}
      \list{\@biblabel{\@arabic\c@enumiv}}%
           {\settowidth\labelwidth{\@biblabel{#1}}%
            \leftmargin\labelwidth
            \advance\leftmargin\labelsep
            \itemsep1pt\parsep0pt\parskip0pt
            \usecounter{enumiv}%
            \let\p@enumiv\@empty
            \renewcommand\theenumiv{\@arabic\c@enumiv}}%
      \sloppy
      \clubpenalty4000
      \@clubpenalty \clubpenalty
      \widowpenalty4000%
      \sfcode`\.\@m}
     {\def\@noitemerr
       {\@latex@warning{Empty `thebibliography' environment}}%
      \endlist}
\makeatother

\begin{document}

\twocolumn[{%
\noindent
\textbf{\Large Memo \#3: AI and the Future of Urban Life} \hfill \textbf{Daniel Hardesty Lewis} \\
\noindent
\textit{Eight Hours: Hybrid Workers, Neighborhood Space, and the Limits of Planning} \hfill April 24, 2026 \\
\textbf{PLAN A6613: AI and the Future of Cities} \hfill Spring 2026
\vspace{3pt}
\hrule
\vspace{8pt}
}]

\section*{I. The Scenario}

\subsection*{Population and Plausibility}

This memo focuses on \textbf{knowledge workers in hybrid employment arrangements} within a dense, transit-oriented metropolitan context: specifically, the New York--Newark--Jersey City metropolitan statistical area. This group is the most plausible early beneficiary of AI-driven time savings for a structurally simple reason: their work is predominantly cognitive and administrative, which is precisely the domain where current generative AI tools deliver measurable productivity gains. McKinsey's 2023 estimate that generative AI could automate 60 to 70 percent of employee time currently spent on ``data processing, data collection, and communicable, predicable physical work'' applies most directly to the information-processing functions that define knowledge labor (McKinsey Global Institute, 2023). This population skews toward workers with at least a bachelor's degree, earning household incomes above the metropolitan median, concentrated in finance, technology, media, legal services, and professional consulting. They are already in partial possession of discretionary time through hybrid arrangements; AI-driven productivity tools represent a further compression of the hours required to deliver equivalent output. The assumption that this group captures eight additional discretionary hours per week is therefore conservative rather than speculative.

\subsection*{How the Time is Reallocated}

This memo assumes that a meaningful share of these freed hours shifts toward \textbf{sustained mid-day and off-peak presence in neighborhood commercial districts and public third places}, at the expense of time previously spent commuting to central business district offices or performing low-value administrative tasks during standard work hours. The specific behavioral shift has three components. First, local café, restaurant, and service spending increases as workers substitute neighborhood commercial districts for Midtown and Downtown lunch corridors. Second, use of neighborhood parks, libraries, branch cultural institutions, and community fitness facilities increases during periods previously occupied by commuting and in-office presence. Third, a smaller but non-trivial share of reallocated time flows toward civic participation: community board attendance, neighborhood association involvement, and local political engagement. The activities workers are shifting \textit{away from} are primarily peak-hour subway trips, desk-bound administrative coordination tasks, and passive in-office presence with no direct productivity function.

\subsection*{Urban Context}

The analysis is grounded in New York City's five boroughs and their immediate transit-connected suburbs. This is a dense, transit-dependent environment where the spatial effects of demand redistribution are unusually legible. Ridership data, commercial district foot traffic, and park utilization all exhibit strong peak-hour concentration. The city's zoning code, public space programming, and transit infrastructure were all designed around an assumption of near-universal 9-to-5 downtown employment. Any systematic redistribution of when and where a large population segment spends its discretionary hours has correspondingly large spatial consequences.

\section*{II. Urban Impacts}

\subsection*{Spatial Demand}

The primary beneficiaries of redistributed time are neighborhood commercial districts that currently experience sharp mid-day troughs in the outer boroughs and in transit-connected neighborhoods that lack significant office employment. Jackson Heights, Astoria, Flatbush, and Park Slope all contain the commercial and public infrastructure to absorb expanded mid-day foot traffic; each has substantial café, restaurant, and retail density whose current utilization curves peak on weekends and dip sharply during weekday working hours. Libraries and branch cultural institutions would see sustained increases in mid-day demand from workers seeking productive, low-cost alternatives to home offices and commercial cafés. New York City's public library system already functions as an informal coworking infrastructure, with branch locations heavily utilized during the pandemic-era remote work surge; a permanent increase in hybrid-worker presence would formalize and intensify this demand pattern.

The spatial losers are concentrated in the central business district. Midtown and Downtown Manhattan's lunch-hour restaurant and retail economy is structurally dependent on the captive audience of workers who are physically present in office towers between 11 AM and 2 PM. A reduction in that captive audience translates directly into revenue pressure for the high-rent, high-turnover food-service businesses that anchor the ground floors of commercial towers. This is not a hypothetical; post-pandemic office attendance data already shows Midtown foot traffic at approximately 72 percent of its pre-pandemic baseline, with average daily attendance at 56 percent of the pre-pandemic rate and the shortfall concentrated on Mondays and Fridays (Partnership for New York City, 2024). AI-accelerated remote and hybrid work deepens an existing trajectory rather than initiating a novel one.

\subsection*{Infrastructure and Services}

The transit system faces a counterintuitive challenge: redistributed demand is not necessarily less demand, but it is \textit{differently distributed} demand that the current system is ill-equipped to serve. The MTA's operating model is optimized for radial, peak-hour flows into Manhattan. AI-driven hybrid work produces a demand pattern that is more distributed across time and more crosstown and intraborough in geography. A worker who previously rode the 7 train into Midtown at 8:30 AM and returned at 6:30 PM now makes more frequent, shorter, midday trips within Queens or between neighborhoods. The MTA currently runs reduced service on crosstown and local routes during off-peak hours, producing a mismatch between where and when the new demand exists and where the system has allocated capacity. Increased mid-day library, park, and café utilization also places pressure on branch library staffing, parks maintenance, and sanitation services that are currently budgeted at off-peak levels. Municipal service delivery was designed around the same peak-hour spatial model as transit; redistributed demand exposes the same structural gap.

The deeper problem is that no existing municipal service system in New York City operates on demand-responsive logic at the neighborhood scale. Branch libraries set their hours based on historical usage data; parks capital budgets are allocated by an annual formula driven by existing park acreage and condition assessments; sanitation routes are fixed to the residential density of the block, not to the actual volume of commercial activity on any given weekday. A population shift of the kind described in this scenario would produce real service gaps, not because the city lacks the resources to address them, but because its administrative systems are not designed to detect and respond to diffuse, gradual, neighborhood-level demand increases. The appropriate planning response is therefore not simply more resources, but a new service-monitoring architecture: granular mid-day foot traffic sensing at branch libraries, parks, and neighborhood commercial corridors, feeding a responsive service-allocation dashboard rather than a fixed annual budget cycle.

\subsection*{Equity}

The equity implications are structurally adverse and should not be minimized. The population segment that captures AI-driven productivity gains is disproportionately white, college-educated, and high-income. The workers whose labor is most threatened by AI automation are disproportionately lower-income, less credentialed, and concentrated in the physical and service sectors that currently generate the foot traffic that sustains neighborhood commercial economies. The city's restaurant delivery workers, retail employees, and building service staff have no equivalent pathway to AI-generated discretionary time; their work is physical, present, and not substitutable by large language models on any near-term horizon. The result is a divergence that is both spatial and economic: hybrid knowledge workers gain time and purchasing power to colonize neighborhood third places, while the workers who staff those spaces face automation of adjacent tasks, wage compression, and intensifying competition for the limited supply of physical-sector employment.

The risk of spatial inequality is also direct. As hybrid worker demand raises the commercial value of neighborhood districts that currently have mixed-income commercial tenants, the upward rent pressure generated by increased foot traffic will accelerate the displacement of lower-margin businesses that serve lower-income residents. The causal chain is the same one this course has examined in the context of algorithmic transit redlining (Mehan, 2026) and AI-driven site feasibility tools (TestFit, 2024): tools that optimize for the preferences of high-income users generate data that computationally justifies capital allocation toward those users, deepening existing geographic inequality. The difference here is that the displacement mechanism is behavioral rather than algorithmic, but the spatial outcome is structurally identical.

\section*{III. Policy and Design Implications}

\subsection*{Land Use and Zoning}

Existing zoning in New York City's residential neighborhoods systematically undersupplies the types of spaces that a population with more discretionary time would most use. The R6 and R7 residential zones that dominate much of Brooklyn, Queens, and the Bronx permit ground-floor retail only along designated commercial overlays, creating long residential blocks with no permitted food service, coworking, or assembly use. As hybrid workers migrate their daytime presence into neighborhoods, they encounter a built environment designed to be passed through rather than inhabited during working hours. The most important land use intervention is a targeted expansion of commercial overlay zoning in transit-accessible residential corridors to permit low-intensity, noise-sensitive commercial uses, specifically cafés, coworking spaces, branch library annexes, and small cultural venues, without requiring full C1 or C2 commercial conversions that would generate ground-floor retail at the scale and rent levels that displace community-serving tenants.

Building typologies present a parallel constraint. The city has permitted the construction of a substantial number of live-work loft conversions and residential towers with shared amenity floors, but these serve the residents of individual buildings rather than the neighborhood. A systematic expansion of the publicly accessible ground-floor amenity model, already demonstrated by the Hudson Yards and Essex Crossing developments, into lower-income neighborhoods through inclusionary development requirements would create publicly accessible coworking, community space, and library infrastructure at the neighborhood scale without requiring the city to fund stand-alone construction.

\subsection*{Economic Development}

The local businesses most likely to benefit from sustained hybrid-worker mid-day presence are neighborhood food and beverage establishments, independent bookstores and cultural retailers, personal service businesses, and fitness and wellness providers in transit-accessible outer-borough neighborhoods. The businesses most at risk are the high-volume, high-rent lunch establishments in the central business district whose economic model depends on the density of captive office workers. City government has an existing toolkit for managing this transition: the Small Business Resilience program, the commercial rent assistance provisions of Local Law 37 (2024), and the BID grant programs administered through the Department of Small Business Services. The appropriate intervention is to redirect a share of these resources toward outer-borough neighborhood commercial districts experiencing increased hybrid-worker demand, specifically to provide rent stabilization support for community-serving tenants who would otherwise be displaced by the commercial rent pressure that increased foot traffic generates.

The risk of unmanaged economic development is that the pattern documented in the wake of the pandemic's initial hybrid-work surge repeats: a small number of high-income neighborhoods capture the bulk of redistributed commercial activity, while the CBD loses at scale and lower-income neighborhoods gain nothing because they lack the commercial infrastructure to capture the demand shift. Active economic development policy must target the undersupplied nodes: neighborhoods with strong transit access, existing residential density, and limited mid-day commercial activity, which are precisely the neighborhoods where rezoning to permit expanded commercial use would generate the greatest equity benefit.

\subsection*{Looking Ahead}

The single most important policy intervention a city should make today to prepare for AI-driven time redistribution is the creation of a \textbf{Neighborhood Time Infrastructure Fund}: a dedicated capital allocation mechanism, housed within the Department of City Planning in partnership with the Department of Cultural Affairs and the public library system, whose explicit mandate is to close the gap between where hybrid workers are migrating their daytime presence and where the public infrastructure to absorb that presence currently exists. The fund would finance branch library expansions and extended hours, publicly accessible coworking annexes in community centers, parks capital improvements that extend usability to year-round and all-weather contexts, and ground-floor commercial activation in underserved residential corridors.

The key stakeholders are the Department of City Planning, the three public library systems, the Parks Department, Borough Presidents with neighborhood commercial constituencies, and local community development financial institutions capable of deploying commercial rent stabilization tools at the neighborhood scale. Successful implementation requires the city to resist two failure modes it has demonstrated in prior AI-adjacent planning contexts. The first is the accountability diffusion documented in the MyCity chatbot's lifecycle: distributing responsibility across enough agencies and vendors that no single office bears accountability for outcomes (Vice, 2026). The second is the algorithmic optimization trap documented in transit AI deployments: using data on existing foot traffic patterns to allocate resources, which systematically underinvests in the under-trafficked neighborhoods that most need infrastructure precisely because their underuse generates the data signals that justify continued underinvestment (Mehan, 2026; AI Now Institute, 2023). The Neighborhood Time Infrastructure Fund must be explicitly counter-cyclical in its allocation logic, directing capital toward demand that does not yet exist in the data but that equitable policy should be designed to generate.\footnote{Full process documentation is available in the project repository at \texttt{memo3/}.}

\begin{thebibliography}{9}
\bibitem{ainow2023} AI Now Institute. (2023). \textit{Algorithmic Redlining and the Calcification of Spatial Inequality}. AI Now Landscape. \url{https://ainowinstitute.org/}
\bibitem{mckinsey2023} McKinsey Global Institute. (2023). \textit{The economic potential of generative AI: The next productivity frontier}. McKinsey \& Company. \url{https://www.mckinsey.com/capabilities/mckinsey-digital/our-insights/the-economic-potential-of-generative-ai}
\bibitem{mehan2026} Mehan, A. (2026). Urban Artificial Intelligence in Mobility Infrastructure. \textit{Contesti Journal}. \url{https://oajournals.fupress.net/index.php/contesti/article/view/16690}
\bibitem{pnyc2024} Partnership for New York City. (2024). \textit{Return to Office Survey: New York City Office Attendance Trends}. \url{https://www.pfnyc.org/}
\bibitem{testfit2024} TestFit. (2024). About TestFit: Our mission, vision, and team. \url{https://www.testfit.io/about}
\bibitem{vice2026} Vice. (2026, February). New York City kills AI chatbot after it kept telling people to break the law. \textit{Vice}. \url{https://www.vice.com/article/nyc-kills-ai-chatbot-mycity/}
\end{thebibliography}

\end{document}
